\section{Week - 10}
\begin{physicsmodule}{Mar 01 - Mar 07}{RoyalBlue}
  \begin{lawbox}
    \subsection*{\sffamily\bfseries Week - 10}
    \begin{enumerate}
    \item Work on completing the results and discussion section particularly for the thickness part.
    \item Also, complete the writing for the hysteresis and curie temperature section for the paper. [Wednesday - whole day]
    \item Go to Rajnandini's Lab one day and spend the time there. [She will tell that]
    \item Work on the netlify deployment and edit that. [Thu-after lunch]
    \item Make the files for the PSSW system calculations and then submit them by the end of the week. [Fri]
\end{enumerate}
  \end{lawbox}
\end{physicsmodule}
%%%%%%%%%%%%%%%%%%%%%%%%%%%%%%%%%%%%%%%%%%%%%%%%%%%%%%%%%%%%%%%%
\subsection*{\underline{Sunday Mar 01, 2026}}
\subsection*{To do Tasks:}
\begin{enumerate}
    \item Make the Agenda's of the meeting in file format.
    \item Do the revision of the bulk and surface damping thing. 
    \item Have family time.
    \item Go for groceries shopping with wife after lunch.
    
\end{enumerate}
\subsection*{Focused Tasks of the day}
\begin{enumerate}[\ding{212}]  
    \item Make the plan of action for the rest of the week daywise and stick to it
\end{enumerate}
%%%%%%%%%%%%%%%%%%%%%%%%%%%%%%%%%%%%%%%%%%%%%%%%%%%%%%%%%%%%%%%%
\subsection*{\underline{Monday Mar 02, 2026}}
\subsection*{To do Tasks:}
\begin{enumerate}
    \item Arrange the papers in the Zotero.
\end{enumerate}
\subsection*{Focused Tasks of the day}
\begin{enumerate}[\ding{212}]  
        \item Arrange the papers in Zotero.
    \item Re-read the thickness dependent measurements of Gilbert damping parameter.
    \item 
\end{enumerate}
\subsection{Important Notes for new papers:}
\begin{enumerate}
    \item The papers which I don't have access and I need for the thickness dependent measured results.
    \begin{enumerate}[(1.)]
        \item \textbf{Magnetic properties of field-annealed FeCo thin films, \textit{JMMM} 320, 739-742 (2008)}
        \begin{enumerate}[\ding{217}]
            \item Large Saturation magnetization, low Coercivity, soft magnetism.
            \item high PMA in $\mathrm{Co_{50}Fe_{50}}$
            \item The dependence of the coercive field with sample thickness is usually hard to predict, and is influenced by a number of factors. An increase of the coercive field on increasing film thickness has been attributed to stresses induced during the deposition in highly magnetostrictive materials.
        \end{enumerate}
        \item \textbf{Non-local detection of spin-polarized electrons at room temperature in $\mathrm{Co_{50}Fe_{50}/GaAs}$ , \textit{Appl. Phys. Lett. } 99, 082108 (2011)}
        \begin{enumerate}[\ding{217}]
            \item This paper is well example of the spin-valve usage of $\mathrm{Co_{50}Fe_{50}}$ and is well for introduction part of the paper.
        \end{enumerate}
        \item \textbf{High-resistivity soft magnetic thin films using CoFe Metal/Native oxide Multilayers (Invided), \textit{Conference on Ferrites}(2004)}
        \begin{enumerate}[\ding{217}]
            \item This paper is well example of the spin-valve usage of $\mathrm{Co_{50}Fe_{50}}$ and is well for introduction part of the paper.
        \end{enumerate}
        \item \textbf{Influence of a Dy overlayer on the precessional dynamics of a ferromagnetic thin film, \textit{Appl. Phys. Lett. } 102, 062418 (2013)}
        \begin{enumerate}[\ding{217}]
            \item Increment in the damping and gyromagnetic ratio of the film while leaving its magnetic anisotropy effectively unchanged.
        \end{enumerate}
        \item \textbf{Influence of a Dy overlayer on the precessional dynamics of a ferromagnetic thin film, \textit{Appl. Phys. Lett. } 102, 062418 (2013)}
        \begin{enumerate}[\ding{217}]
            \item Increment in the damping and gyromagnetic ratio of the film while leaving its magnetic anisotropy effectively unchanged.
        \end{enumerate}
        \item \textbf{Spin pumping in Ferromagnet-Topological Insulator-Ferromagnet Heterostructures, \textit{Scientific Reports} 5, 7907 (2015)}
        \begin{enumerate}[\ding{217}]
            \item Spin pumping study for the TI-ferromagnet interface, investigating spin transfer dynamics in a spin-valve like structure.
            \item Dynamic coupling is limited as per the layer-resolved measurements.
        \end{enumerate}
        \item \textbf{Spin pumping in magnetic trilayer structures with an MgO barrier, \textit{Scientific Reports} 6, 35582 (2016)}
        \begin{enumerate}[\ding{217}]
            \item There is an increase in the interlayer exchange coupling when the MgO's thickness is increased from 1 nm to 4 nm in the trilayer system.
            \item Gilbert damping is found to be independent of the MgO thickness, suggesting the suppression of spin-pumping.
            \item There are some spins that are pumped by the precessing magnetization of the CoFe are absorbed. The lower damping in the case of the thinnest MgO layer could be in part due to the strong static exchange coupling observed.
        \end{enumerate}
        \item \textbf{Engineering strong magnon-phonon coupling in single CoFe nanomagnets, \textit{J. Appl. Phys.} 137, 123904 (2025)} - Important paper
        \begin{enumerate}[\ding{217}]
            \item The strong coupling between magnonic and phonic modes in both polycrystalline and single-crystalline CoFe nanomagnets is observed using TRMOKE.
            \item The enhanced coupling strength in the single-crystalline CoFe nanomagnet is attributed to the higher magnetorestriction in CoFe systems and lower intrinsic magnetic damping.
            \item The relative values for the $\mathrm{\alpha_{eff}}$ was observed for different frequencies.
        \end{enumerate}
        \item \textbf{Interfacial Tuning of Anisotropic Gilbert Damping, \textit{Phys. Rev. Lett.} 130, 046704 (2023)} 
        \begin{enumerate}[\ding{217}]
            \item The sign reversal of the anisotropic damping as well as the anisotropic effective spin mixing conductance in ultrathin Fe films grown on GaAs(001), when changing the thickness of the Fe.
            \item \textcolor{blue}{Nice Paper - Read later in peace.}
        \end{enumerate}
        \item \textbf{Gilbert damping Parameter in MgO-based Magnetic Tunnel Junctions from First principles, \textit{Phys. Rev. Applied} 7, 034004 (2017)} - \textcolor{blue}{Important Paper}  
        \begin{enumerate}[\ding{217}]
            \item The calculated dependence of Gilbert Damping on the thickness of the MgO capping layer is consistent with experiment and indicates that the deceases in $\alpha$ with increasing thickness of the MgO capping layer is caused by suppression of spin pumping.
        \end{enumerate}
    \end{enumerate}
\end{enumerate}
% %%%%%%%%%%%%%%%%%%%%%%%%%%%%%%%%%%%%%%%%%%%%%%%%%%%%%%%%%%%%%%%%%
\subsection*{\underline{Tuesday Mar 03, 2026}}
\subsection*{To do Tasks:}
\begin{enumerate}
    \item Rearrange Zotero and then start writing the thickness part. \cmark
    \item Additionally, choose the thickness and try to create the simulation folder. \cmark
    \item Add Sanjay bhaiya's site to personal repo and activate the required theme later in the day. \xmark
\end{enumerate}
\subsection*{Focused Tasks of the day}
\begin{enumerate}[\ding{212}]  
    \item Complete the thickness-dependent work for the paper.
\end{enumerate}
\subsection*{Notes:}
\begin{enumerate}
    \item Need to arrange the text in the Word tomorrow after completing the surface \& bulk part writing, and also the equilibrium dynamics related things.
\end{enumerate}
% %%%%%%%%%%%%%%%%%%%%%%%%%%%%%%%%%%%%%%%%%%%%%%%%%%%%%%%%%%%%%%%%%
\subsection*{\underline{Wednesday Mar 04, 2026} }
\subsection*{To do Tasks:}
\begin{enumerate}
    \item Complete the remaining writing for Surface/Bulk damping and other equilibrium details. [\textcolor{blue}{partial}]
    \item Complete the personal website up to date. \cmark
    \item Attend the meeting with GSA at 4:00 pm. \cmark
    \item Add the online page of the GRS-2026 talk with an image and post the link also on LinkedIn.
    \item Make the Agenda File and share that in the email.
    \item Write some part of the \textbf{Force-Note} before bed.
    \item Go to the temple in the morning if possible, as it is Holi. \xmark
\end{enumerate}
\subsection*{Focused Tasks of the day}
\begin{enumerate}[\ding{212}]  
    \item Work on writing in the morning.
    \item Complete the website and make it up to date. \cmark
\end{enumerate}
% \subsection*{Notes:}
% \begin{enumerate}
%     \item 
% \end{enumerate}
% % %%%%%%%%%%%%%%%%%%%%%%%%%%%%%%%%%%%%%%%%%%%%%%%%%%%%%%%%%%%%%%%%%
\subsection*{\underline{Thursday Mar 05, 2026}}
\subsection*{To do Tasks:}
\begin{enumerate}
    \item Make the meeting agenda for Notes for Physics Meeting. \cmark
    \item Start the writing from the morning for the paper. \cmark
    \item Design the LinkedIn post for the GRS and also a news page with a photo of Aidan and me on the website after completing the surface and bulk damping writing.
    \item Work on writing the note for \textbf{Force}.
\end{enumerate}
\subsection*{Focused Tasks of the day}
\begin{enumerate}[\ding{212}]  
    \item Complete the Agenda for Notes for Physics Meeting. \cmark
    \item Write the surface and bulk damping-related part of the paper. \cmark 
    \item Go to the gym later in the day after lunch.
    \item Work on writing the details of the system and parameters taken with proper source, and also start the hysteresis and Curie temperature part before bed today.
\end{enumerate}
\subsection*{Notes:}
\begin{enumerate}
    \item 
\end{enumerate}
% %%%%%%%%%%%%%%%%%%%%%%%%%%%%%%%%%%%%%%%%%%%%%%%%%%%%%%%%%%%%%%%%%
\subsection*{\underline{Friday Mar 06, 2026}}
\subsection*{To do Tasks:}
\begin{enumerate}
    \item Add Suwash's website to the website.
    \item 
\end{enumerate}
\subsection*{Focused Tasks of the day}
\begin{enumerate}[\ding{212}]  
    \item 
\end{enumerate}
\subsection*{Notes:}
\begin{enumerate}
    \item 
\end{enumerate}
% % %%%%%%%%%%%%%%%%%%%%%%%%%%%%%%%%%%%%%%%%%%%%%%%%%%%%%%%%%%%%%%%%
\subsection*{\underline{Saturday Mar 07, 2026} - My Birthday }
\subsection*{To do Tasks:}
\begin{enumerate}
    \item Enjoy the day and be happy.
\end{enumerate}
% \subsection*{Focused Tasks of the day}
% \begin{enumerate}[\ding{212}]  
%     \item 
% \end{enumerate}
\subsection*{Notes:}
\begin{enumerate}
    \item Went to Sri Laxmi Temple.
    \item Later went to Prudential Center for \textbf{View Boston} thing and from there we went to The Cheesecake Factory for dinner.
    \item A day well spent with my wife.
\end{enumerate}
\stardivider
